%%%%%%%%%%%%%%%%%%%%%%%%%%%%%%%%%%%%%%%%%%%%%%%%%%%%%%%%%%%%%%%%%%%%%%
% Lớp 9 | Chương VIII. Xác suất của biến cố trong một số mô hình xác suất đơn giản | Xác suất của biến cố liên quan tới phép thử
% Dạng 9D8?26-3: Bài toán xác suất thực tế
%%%%%%%%%%%%%%%%%%%%%%%%%%%%%%%%%%%%%%%%%%%%%%%%%%%%%%%%%%%%%%%%%%%%%%


% nguồn: 2026-11-14_09_de_tuyen_sinh_10_tphcm_2025_2026 (nhập 14/11/2026)
\begin{bt}%[9D8H26-3]%[ma:b00cb57d]
% nguon: Đề thi chính thức Tuyển sinh 10 TP.HCM năm học 2025-2026 – Bài 3
Biểu đồ tròn cho biết tỉ lệ về số lượng các loại bảo hiểm đã bán được trong tháng 4/2025 của một công ty. Biết rằng trong tháng này, công ty đã bán được $300$ gói bảo hiểm các loại cho $300$ khách hàng khác nhau:
\begin{itemize}
	\item Loại bảo hiểm A: $25\%$
	\item Loại bảo hiểm B: $15\%$
	\item Loại bảo hiểm C: $33\%$
	\item Loại bảo hiểm D: $27\%$
\end{itemize}
\begin{enumerate}[a)]
	\item Tính số lượng cụ thể của mỗi loại bảo hiểm mà công ty đã bán được trong tháng 4/2025?
	\item Bộ phận chăm sóc khách hàng chọn ngẫu nhiên một khách hàng đã mua bảo hiểm của công ty trong tháng 4/2025 để khảo sát. Tính xác suất của biến cố: ``Khách hàng được chọn không mua loại bảo hiểm B''.
\end{enumerate}
\loigiai{
\begin{enumerate}[a)]
	\item Số lượng cụ thể của từng loại bảo hiểm công ty đã bán được trong tháng 4/2025 là:
	\begin{itemize}
		\item Số lượng gói bảo hiểm loại A: $300 \cdot 25\% = 75$ (gói).
		\item Số lượng gói bảo hiểm loại B: $300 \cdot 15\% = 45$ (gói).
		\item Số lượng gói bảo hiểm loại C: $300 \cdot 33\% = 99$ (gói).
		\item Số lượng gói bảo hiểm loại D: $300 \cdot 27\% = 81$ (gói).
	\end{itemize}
	\item Số các kết quả có thể xảy ra khi chọn ngẫu nhiên 1 khách hàng là $n(\Omega) = 300$.\\
	Gọi biến cố $E$: ``Khách hàng được chọn không mua loại bảo hiểm B''.\\
	Khi đó khách hàng được chọn mua loại bảo hiểm A, C hoặc D. Số kết quả thuận lợi cho biến cố $E$ là:
	$n(E) = 75 + 99 + 81 = 255$ (khách hàng).\\
	Xác suất của biến cố $E$ là:
	$P(E) = \dfrac{n(E)}{n(\Omega)} = \dfrac{255}{300} = \dfrac{17}{20} = 0{,}85$.\\
	Vậy xác suất cần tìm là $\dfrac{17}{20}$.
\end{enumerate}
}
\end{bt}
